\begin{textsample}{POS Dim 1 – 2010 – Score 28.00 – 2010/t002801.txt}  \label{ex:f1_pos_028}
I would like to share with you this morning some stories about \textbf{the} ocean through my work as a still photographer for National Geographic magazine.
I guess I became an \textbf{underwater} photographer and a photojournalist because I fell in love with \textbf{the} \textbf{sea} as a child.
And I wanted to tell stories about all \textbf{the} amazing things I was seeing \textbf{underwater}, incredible wildlife and interesting behaviors.
And after even 30 years of doing this, after 30 years of exploring \textbf{the} ocean, I never cease to be amazed at \textbf{the} extraordinary encounters that I have while I ' m at \textbf{sea}.
But more and more frequently these days I ' m seeing terrible things \textbf{underwater} as well, things that I don ' t think most people realize.
And I ' ve been compelled to turn my camera towards these issues to tell a more complete story.
I want people to see what ' s happening \textbf{underwater}, both \textbf{the} horror and \textbf{the} magic. \textbf{The} first story that I did for National Geographic, where I recognized \textbf{the} ability to include environmental issues within a natural history coverage, was a story I proposed on harp \textbf{seals}. \textbf{The} story I wanted to do initially was just a small focus to look at \textbf{the} few weeks each year where these animals migrate down from \textbf{the} Canadian arctic to \textbf{the} Gulf of St.
Lawrence in Canada to engage in courtship, mating and to have their pups.
And all of this is played out against \textbf{the} backdrop of transient pack ice that moves with wind and tide.
And because I ' m an \textbf{underwater} photographer, I wanted to do this story from both above and below, to make pictures like this that show one of these little pups making its very first swim in \textbf{the} icy 29-degree water.
But as I got more involved in \textbf{the} story, I realized that there were two big environmental issues I couldn ' t ignore. \textbf{The} first was that these animals continue to be hunted, killed with hakapiks at about eight, 15 days old.
It actually is \textbf{the} largest marine mammal slaughter on \textbf{the} \textbf{planet}, with hundreds of thousands of these \textbf{seals} being killed every year.
But as disturbing as that is, I think \textbf{the} bigger problem for harp \textbf{seals} is \textbf{the} loss of \textbf{sea} ice due to global warming.
This is an aerial picture that I made that shows \textbf{the} Gulf of St.
Lawrence during harp \textbf{seal} season.
And even though we see a lot of ice in this picture, there ' s a lot of water as well, which wasn ' t there historically.
And \textbf{the} ice that is there is quite thin. \textbf{The} problem is that these pups need a stable platform of solid ice in order to nurse from their moms.
They only need 12 days from \textbf{the} moment they ' re born until they ' re on their own.
But if they don ' t get 12 days, they can fall into \textbf{the} ocean and die.
This is a photo that I made showing one of these pups that ' s only about five or seven days old — still has a little bit of \textbf{the} umbilical cord on its belly — that has fallen in because of \textbf{the} thin ice, and \textbf{the} mother is frantically trying to push it up to breathe and to get it back to stable purchase.
This problem has continued to grow each year since I was there.
I read that last year \textbf{the} pup mortality rate was 100 percent in parts of \textbf{the} Gulf of St.
Lawrence.
So, clearly, this species has a lot of problems going forward.
This ended up becoming a cover story at National Geographic.
And it received quite a bit of attention.
And with that, I saw \textbf{the} potential to begin doing other stories about ocean problems.
So I proposed a story on \textbf{the} global \textbf{fish} crisis, in part because I had personally witnessed a lot of degradation in \textbf{the} ocean over \textbf{the} last 30 years, but also because I read a scientific paper that stated that 90 percent of \textbf{the} big \textbf{fish} in \textbf{the} ocean have \textbf{disappeared} in \textbf{the} last 50 or 60 years.
These are \textbf{the} tuna, \textbf{the} billfish and \textbf{the} sharks.
When I read that, I was blown away by those numbers.
I thought this was going to be headline news in every media outlet, but it really wasn ' t, so I wanted to do a story that was a very different kind of \textbf{underwater} story.
I wanted it to be more like war photography, where I was making harder-hitting pictures that showed readers what was happening to marine wildlife around \textbf{the} \textbf{planet}. \textbf{The} first component of \textbf{the} story that I thought was essential, however, was to give readers a sense of appreciation for \textbf{the} ocean animals that they were eating.
You know, I think people go into a restaurant, and somebody orders a steak, and we all know where steak comes from, and somebody orders a chicken, and we know what a chicken is, but when they ' re eating bluefin sushi, do they have any sense of \textbf{the} magnificent animal that they ' re consuming?
These are \textbf{the} \textbf{lions} and \textbf{tigers} of \textbf{the} \textbf{sea}.
In reality, these animals have no terrestrial counterpart ; they ' re unique in \textbf{the} world.
These are animals that can practically swim from \textbf{the} equator to \textbf{the} poles and can crisscross entire oceans in \textbf{the} course of a year.
If we weren ' t so efficient at catching them, because they grow their entire life, would have 30-year-old bluefin out there that weigh a ton.
But \textbf{the} truth is we ' re way too efficient at catching them, and their stocks have collapsed worldwide.
This is \textbf{the} daily auction at \textbf{the} Tsukiji Fish Market that I photographed a couple years ago.
And every single day these tuna, bluefin like this, are \textbf{stacked} up like cordwood, just warehouse after warehouse.
As I wandered around and made these pictures, it sort of \textbf{occurred} to me that \textbf{the} ocean ' s not a grocery store, you know.
We can ' t keep taking without expecting serious consequences as a result.
I also, with \textbf{the} story, wanted to show readers how \textbf{fish} are caught, some of \textbf{the} methods that are used to catch \textbf{fish}, like a bottom trawler, which is one of \textbf{the} most common methods in \textbf{the} world.
This was a small net that was being used in Mexico to catch \textbf{shrimp}, but \textbf{the} way it works is essentially \textbf{the} same everywhere in \textbf{the} world.
You have a large net in \textbf{the} middle with two steel doors on either end.
And as this assembly is towed through \textbf{the} water, \textbf{the} doors meet resistance with \textbf{the} ocean, and it opens \textbf{the} mouth of \textbf{the} net, and they place floats at \textbf{the} top and a lead line on \textbf{the} bottom.
And this just drags over \textbf{the} bottom, in this case to catch \textbf{shrimp}.
But as you can imagine, it ' s catching everything else in its path as well.
And it ' s destroying that precious benthic community on \textbf{the} bottom, things like sponges and corals, that critical \textbf{habitat} for other animals.
This photograph I made of \textbf{the} fisherman holding \textbf{the} \textbf{shrimp} that he caught after towing his nets for one hour.
So he had a handful of \textbf{shrimp}, maybe seven or eight \textbf{shrimp}, and all those other animals on \textbf{the} deck of \textbf{the} boat are bycatch.
These are animals that died in \textbf{the} process, but have no commercial value.
So this is \textbf{the} true cost of a \textbf{shrimp} dinner, maybe seven or eight \textbf{shrimp} and 10 \textbf{pounds} of other animals that had to die in \textbf{the} process.
And to make that point even more visual, I swam under \textbf{the} \textbf{shrimp} boat and made this picture of \textbf{the} guy shoveling this bycatch into \textbf{the} \textbf{sea} as trash and photographed this cascade of death, you know, animals like guitarfish, bat rays, flounder, pufferfish, that only an hour before, were on \textbf{the} bottom of \textbf{the} ocean, alive, but now being thrown back as trash.
I also wanted to focus on \textbf{the} shark \textbf{fishing} industry because, currently on \textbf{planet} Earth, we ' re killing over 100 million sharks every single year.
But before I went out to photograph this component, I sort of wrestled with \textbf{the} notion of how do you make a picture of a dead shark that will resonate with readers You know, I think there ' s still a lot of people out there who think \textbf{the} only good shark is a dead shark.
But this one morning I jumped in and found this thresher that had just recently died in \textbf{the} gill net.
And with its huge pectoral \textbf{fins} and eyes still very visible, it struck me as sort of a crucifixion, if you will.
This ended up being \textbf{the} lead picture in \textbf{the} global fishery story in National Geographic.
And I hope that it helped readers to take notice of this problem of 100 million sharks.
And because I love sharks — I ' m somewhat obsessed with sharks — I wanted to do another, more celebratory, story about sharks, as a way of talking about \textbf{the} need for shark conservation.
So I went to \textbf{the} Bahamas because there ' re very few places in \textbf{the} world where sharks are doing well these days, but \textbf{the} Bahamas seem to be a place where stocks were reasonably healthy, largely due to \textbf{the} fact that \textbf{the} government there had outlawed longlining several years ago.
And I wanted to show several species that we hadn ' t shown much in \textbf{the} magazine and worked in a number of locations.
One of \textbf{the} locations was this place called Tiger Beach, in \textbf{the} northern Bahamas where \textbf{tiger} sharks aggregate in shallow water.
This is a low-altitude photograph that I made showing our dive boat with about a dozen of these big old \textbf{tiger} sharks sort of just swimming around behind.
But \textbf{the} one thing I definitely didn ' t want to do with this coverage was to continue to portray sharks as something like monsters.
I didn ' t want them to be overly threatening or scary.
And with this photograph of a beautiful 15-feet, probably 14-feet, I guess, female \textbf{tiger} shark, I sort of think I got to that goal, where she was swimming with these little barjacks off her nose, and my strobe created a shadow on her face.
And I think it ' s a gentler picture, a little less threatening, a little more respectful of \textbf{the} species.
I also \textbf{searched} on this story for \textbf{the} elusive great hammerhead, an animal that really hadn ' t been photographed much until maybe about seven or 10 years ago.
It ' s a very solitary creature.
But this is an animal that ' s considered data deficient by science in both Florida and in \textbf{the} Bahamas.
You know, we know almost nothing about them.
We don ' t know where they migrate to or from, where they \textbf{mate}, where they have their pups, and yet, hammerhead populations in \textbf{the} Atlantic have declined about 80 percent in \textbf{the} last 20 to 30 years.
You know, we ' re losing them faster than we can possibly find them.
This is \textbf{the} oceanic whitetip shark, an animal that is considered \textbf{the} fourth most dangerous species, if you pay attention to such lists.
But it ' s an animal that ' s about 98 percent in decline throughout most of its range.
Because this is a pelagic animal and it lives out in \textbf{the} deeper water, and because we weren ' t working on \textbf{the} bottom, I brought along a shark cage here, and my friend, shark biologist Wes Pratt is inside \textbf{the} cage.
You ' ll see that \textbf{the} photographer, of course, was not inside \textbf{the} cage here, so clearly \textbf{the} biologist is a little smarter than \textbf{the} photographer I guess.
And lastly with this story, I also wanted to focus on baby sharks, shark nurseries.
And I went to \textbf{the} \textbf{island} of Bimini, in \textbf{the} Bahamas, to work with lemon shark pups.
This is a photo of a lemon shark pup, and it shows these animals where they live for \textbf{the} first two to three years of their lives in these protective mangroves.
This is a very sort of un-shark-like photograph.
It ' s not what you typically might think of as a shark picture.
But, you know, here we see a shark that ' s maybe 10 or 11 \textbf{inches} long swimming in about a foot of water.
But this is crucial \textbf{habitat} and it ' s where they spend \textbf{the} first two, three years of their lives, until they ' re big enough to go out on \textbf{the} rest of \textbf{the} reef.
After I left Bimini, I actually learned that this \textbf{habitat} was being bulldozed to create a new \textbf{golf} course and resort.
And other recent stories have looked at single, flagship species, if you will, that are at risk in \textbf{the} ocean as a way of talking about other threats.
One such story I did documented \textbf{the} leatherback \textbf{sea} \textbf{turtle}.
This is \textbf{the} largest, widest- ranging, deepest-diving and oldest of all \textbf{turtle} species.
Here we see a female crawling out of \textbf{the} ocean under moonlight on \textbf{the} \textbf{island} of Trinidad.
These are animals whose lineage dates back about 100 million years.
And there was a time in their lifespan where they were coming out of \textbf{the} water to nest and saw Tyrannosaurus rex running by.
And today, they crawl out and see condominiums.
But despite this amazing longevity, they ' re now considered critically endangered.
In \textbf{the} Pacific, where I made this photograph, their stocks have declined about 90 percent in \textbf{the} last 15 years.
This is a photograph that shows a hatchling about to taste saltwater for \textbf{the} very first time beginning this long and perilous journey.
Only one in a thousand leatherback hatchlings will reach maturity.
But that ' s due to natural predators like vultures that pick them off on a beach or predatory \textbf{fish} that are waiting offshore.
Nature has learned to compensate with that, and females have multiple clutches of eggs to overcome those odds.
But what they can ' t deal with is anthropogenic stresses, human things, like this picture that shows a leatherback caught at night in a gill net.
I actually jumped in and photographed this, and with \textbf{the} fisherman ' s permission, I cut \textbf{the} \textbf{turtle} out, and it was able to swim free.
But, you know, thousands of other leatherbacks each year are not so fortunate, and \textbf{the} species ' future is in great danger.
Another charismatic megafauna species that I worked with is \textbf{the} story I did on \textbf{the} right \textbf{whale}.
And essentially, \textbf{the} story is this with right \textbf{whales}, that about a million years ago, there was one species of right \textbf{whale} on \textbf{the} \textbf{planet}, but as land masses moved around and oceans became isolated, \textbf{the} species sort of separated, and today we have essentially two distinct stocks.
We have \textbf{the} Southern right \textbf{whale} that we see here and \textbf{the} North Atlantic right \textbf{whale} that we see here with a mom and calf off \textbf{the} coast of Florida.
Now, both species were hunted to \textbf{the} brink of extinction by \textbf{the} early whalers, but \textbf{the} Southern right \textbf{whales} have rebounded a lot better because they ' re located in places farther away from human activity. \textbf{The} North Atlantic right \textbf{whale} is listed as \textbf{the} most endangered species on \textbf{the} \textbf{planet} today because they are urban \textbf{whales} ; they live along \textbf{the} east coast of North America, United States and Canada, and they have to deal with all these urban ills.
This photo shows an animal popping its head out at sunset off \textbf{the} coast of Florida.
You can see \textbf{the} \textbf{coal} burning plant in \textbf{the} background.
They have to deal with things like toxins and pharmaceuticals that are flushed out into \textbf{the} ocean, and maybe even affecting their \textbf{reproduction}.
They also get entangled in \textbf{fishing} gear.
This is a picture that shows \textbf{the} tail of a right \textbf{whale}.
And those white markings are not natural markings.
These are entanglement scars. 72 percent of \textbf{the} population has such scars, but most don ' t shed \textbf{the} gear, things like lobster traps and \textbf{crab} pots.
They hold on to them, and it eventually kills them.
And \textbf{the} other problem is they get hit by ships.
And this was an animal that was struck by a ship in Nova Scotia, Canada being towed in, where they did a necropsy to confirm \textbf{the} cause of death, which was indeed a ship strike.
So all of these ills are \textbf{stacking} up against these animals and keeping their numbers very low.
And to draw a contrast with that beleaguered North Atlantic population, I went to a new pristine population of Southern right \textbf{whales} that had only been discovered about 10 years ago in \textbf{the} sub-Antarctic of New Zealand, a place called \textbf{the} Auckland Islands.
I went down there in \textbf{the} winter time.
And these are animals that had never seen humans before, and I was one of \textbf{the} first people they probably had ever seen.
And I got in \textbf{the} water with them, and I was amazed at how curious they were.
This photograph shows my assistant standing on \textbf{the} bottom at about 70 feet and one of these amazingly beautiful, 45-foot, 70- ton \textbf{whales}, like a city bus just swimming up, you know.
They were in perfect condition, very \textbf{fat} and healthy, robust, no entanglement scars, \textbf{the} way they ' re supposed to look.
You know, I read that \textbf{the} pilgrims, when they landed at Plymouth Rock in Massachusetts in 1620, wrote that you could walk across Cape Cod Bay on \textbf{the} backs of right \textbf{whales}.
And we can ' t go back and see that today, but maybe we can preserve what we have left.
And I wanted to close this program with a story of hope, a story I did on marine reserves as sort of a solution to \textbf{the} problem of overfishing, \textbf{the} global \textbf{fish} crisis story.
I settled on working in \textbf{the} country of New Zealand because New Zealand was rather progressive, and is rather progressive in terms of protecting their ocean.
And I really wanted this story to be about three things : I wanted it to be about abundance, about diversity and about resilience.
And one of \textbf{the} first places I worked was a reserve called Goat Island in Leigh of New Zealand.
What \textbf{the} scientists there told me was that when protected this first marine reserve in 1975, they hoped and expected that certain things might happen.
For example, they hoped that certain species of \textbf{fish} like \textbf{the} New Zealand snapper would return because they had been \textbf{fished} to \textbf{the} brink of commercial extinction.
And they did come back.
What they couldn ' t predict was that other things would happen.
For example, these \textbf{fish} predate on \textbf{sea} urchins, and when \textbf{the} \textbf{fish} were all gone, all anyone ever saw \textbf{underwater} was just acres and acres of \textbf{sea} urchins.
But when \textbf{the} \textbf{fish} came back and began predating and controlling \textbf{the} urchin population, low and behold, kelp forests emerged in shallow water.
And that ' s because \textbf{the} urchins eat kelp.
So when \textbf{the} \textbf{fish} control \textbf{the} urchin population, \textbf{the} ocean was restored to its natural equilibrium.
You know, this is probably how \textbf{the} ocean looked here one or 200 years ago, but nobody was around to tell us.
I worked in other parts of New Zealand as well, in beautiful, fragile, protected areas like in Fiordland, where this \textbf{sea} pen colony was found.
Little \textbf{blue} cod swimming in for a dash of color.
In \textbf{the} northern part of New Zealand, I dove in \textbf{the} \textbf{blue} water, where \textbf{the} water ' s a little warmer, and photographed animals like this giant sting ray swimming through an \textbf{underwater} canyon.
Every part of \textbf{the} ecosystem in this place seems very healthy, from tiny, little animals like a nudibrank crawling over encrusting sponge or a leatherjacket that is a very important animal in this ecosystem because it grazes on \textbf{the} bottom and allows new life to take hold.
And I wanted to finish with this photograph, a picture I made on a very stormy day in New Zealand when I just laid on \textbf{the} bottom amidst a school of \textbf{fish} swirling around me.
And I was in a place that had only been protected about 20 years ago.
And I talked to divers that had been diving there for many years, and they said that \textbf{the} marine life was better here today than it was in \textbf{the} 1960s.
And that ' s because it ' s been protected, that it has come back.
So I think \textbf{the} message is clear. \textbf{The} ocean is, indeed, resilient and tolerant to a point, but we must be good custodians.
I became an \textbf{underwater} photographer because I fell in love with \textbf{the} \textbf{sea}, and I make pictures of it today because I want to protect it, and I don ' t think it ' s too late.
Thank you very much.

% matched lemmas: blue, coal, crab, disappear, fat, fin, fish, fishing, golf, habitat, inch, island, lion, mate, occur, planet, pound, reproduction, sea, seal, search, shrimp, stack, the, tiger, turtle, underwater, whale
\end{textsample}
